\documentclass[12pt]{article}
\usepackage[utf8]{inputenc}

\usepackage{lmodern}

\usepackage{enumitem}
\usepackage[margin=2cm]{geometry}

\usepackage{amsmath, amsfonts, amssymb}
\usepackage{graphicx}
%\usepackage{subfigure}
\usepackage{tikz}
\usepackage{pgfplots}
\usepackage{multicol}

\usepackage{comment}
\usepackage{url}
\usepackage{calc}
\usepackage{subcaption}
\usepackage[indent=0pt]{parskip}
\usepackage{animate}

\usepackage{array}
\usepackage{blkarray,booktabs, bigstrut}
\usepackage{bigints}

\pgfplotsset{compat=1.16}

% MATH commands
\newcommand{\ga}{\left\langle}
\newcommand{\da}{\right\rangle}
\newcommand{\oa}{\left\lbrace}
\newcommand{\fa}{\right\rbrace}
\newcommand{\oc}{\left[}
\newcommand{\fc}{\right]}
\newcommand{\op}{\left(}
\newcommand{\fp}{\right)}

\newcommand{\bi}{\mathbf{i}}
\newcommand{\bj}{\mathbf{j}}
\newcommand{\bk}{\mathbf{k}}
\newcommand{\bF}{\mathbf{F}}

\newcommand{\mR}{\mathbb{R}}

\newcommand{\ra}{\rightarrow}
\newcommand{\Ra}{\Rightarrow}

\newcommand{\sech}{\mathrm{sech}\,}
\newcommand{\csch}{\mathrm{csch}\,}
\newcommand{\curl}{\mathrm{curl}\,}
\newcommand{\dive}{\mathrm{div}\,}

\newcommand{\ve}{\varepsilon}
\newcommand{\spc}{\vspace*{0.5cm}}

\DeclareMathOperator{\Ran}{Ran}
\DeclareMathOperator{\Dom}{Dom}

\newcommand{\exo}[1]{\noindent\textcolor{red}{\fbox{\textbf{Problem {#1}}}\hrulefill}\\}
\newcommand{\qu}[4]{\noindent\textcolor{#4}{\fbox{\textbf{Section {#1} | Problem {#2}}} \hrulefill{{\fbox{\textbf{{#3} Points}}}}\\}}

\newcommand{\semester}{Spring 2023}

\newcommand{\CVup}{%
\begin{tikzpicture}
\draw[black, <->, >=latex] (-0.33, 0.5) .. controls (-0.125, 0) and (0.125, 0) .. (0.33, 0.5);
\end{tikzpicture}}

\newcommand{\CVupInc}{%
\begin{tikzpicture}
\draw[black, ->, >=latex] (0,0) .. controls (0.2, 0) and (0.4, 0.2) .. (0.5, 0.5);
\end{tikzpicture}}

\newcommand{\CVupDec}{%
\begin{tikzpicture}[rotate=270]
\draw[black, ->, >=latex] (0,0) .. controls (0.2, 0) and (0.4, 0.2) .. (0.5, 0.5);
\end{tikzpicture}}

\newcommand{\CVdown}{%
\begin{tikzpicture}
\draw[black, <->, >=latex] (-0.33, -0.5) .. controls (-0.125, 0) and (0.125, 0) .. (0.33, -0.5);
\end{tikzpicture}}

\newcommand{\CVdownInc}{%
\begin{tikzpicture}
\draw[black, ->, >=latex] (-0.5, -0.5) .. controls (-0.5, -0.3) and (-0.5, -0.1) .. (0,0);
\end{tikzpicture}}

\newcommand{\CVdownDec}{%
\begin{tikzpicture}[rotate=-90]
\draw[black, ->, >=latex] (-0.5, -0.5) .. controls (-0.5, -0.3) and (-0.5, -0.1) .. (0,0);
\end{tikzpicture}}

\begin{document}
	\noindent \hrulefill \\
	MATH-241 \hfill Pierre-Olivier Paris{\'e}\\
	Solutions Section 2-9 \hfill \semester \\\vspace*{-1cm}
	
	\noindent\hrulefill
	
	\spc

	\exo{4}
	\\
	The linearization is given by
		\begin{align*}
		L (x) = f(3) + f'(3) (x - 3)  .
		\end{align*}
	We have 
		\begin{align*}
		f' (x) = \frac{-2x}{(x^2 - 5)^{3/2}} .
		\end{align*}
	Therefore, we have $f'(3) = \frac{-6}{8} = -\frac{3}{4}$. We also have $f(3) = 1$. Therefore, the linearization is
		\begin{align*}
		L (z) = 1 -\frac{3}{4} (x - 3)
		\end{align*}	
		
	\spc
	
	\exo{24}
	\\
	We see that
		\begin{align*}
		4.002 = 4 + 0.002 .
		\end{align*}		
	Therefore, the value $0.002$ will by my $x$ in the linearization and suggest
		\begin{align*}
		f(x) = \frac{1}{x + 4} .
		\end{align*}
	We see that $f(0.002) = 1/4.002$. Therefore, a linear approximation of $f$ around $x = 0$ will be useful to approximation $1/4.002$. We have
		\begin{align*}
		f'(x) = - \frac{1}{(x + 4)^2} \quad \Ra \quad f' (0) = - \frac{1}{16} .
		\end{align*}
	Since $f(0) = 1/4$, we have
		\begin{align*}
		L (x) = 1/4 - \frac{x}{16} .
		\end{align*}
	Using the linearization of $f$, we find that
		\begin{align*}
		\frac{1}{4.002} = f(0.002) \approx L(0.002) = 0.25 - \frac{0.002}{16} = 0.25 - 0.000125 = 0.249875 .
		\end{align*}
	Therefore, we have
		\begin{align*}
		\frac{1}{4.002} \approx 0.249875 .
		\end{align*}
		
	
\end{document}